% SEGMENT OLP-0570-S01
\documentclass[../../../include/open-logic-section]{subfiles}

\begin{document}

\olfileid{sth}{replacement}{absinf}
\olsection{மாற்றீடும் ``முற்றுமுடிவிலியும்''}

இப்போது \limofsize{} ஐப் பின்னால் விட்டுவிட்டு, கணங்கள் படிநிலைகளில்
உருவாகின்றன என்ற கருத்தைத் தீவிரமாக எடுத்துக்கொள்ளும், மாற்றீட்டை
நியாயப்படுத்துவதற்கான வேறொரு (உள்ளார்ந்த) முயற்சிக் குடும்பத்தை ஆராய்வோம்.

படிநிலைமுறைச் செயல்முறையை முதலில் வரைந்தபோது, ஒவ்வொரு படிநிலையிலும் என்ன
நிகழ்கிறது என்பதை விளக்கும் சில கொள்கைகளை வழங்கினோம். அவை
\stageshier, \stagesord, மற்றும் \stagesacc. பின்னர் படிநிலைகளின்
எண்ணிக்கையைப் பற்றி ஏதோ ஒன்றைச் சொல்லும் சில கொள்கைகளைச் சேர்த்தோம்:
கணம் உருவாக்கும் செயல்முறை ஒருபோதும் முடிவதில்லை என்று \stagessucc{}
கூறியது; அந்தச் செயல்முறை ஒரு முடிவிலியாம் படிநிலையைக் கடந்து செல்கிறது
என்று \stagesinf{} கூறியது.

இப்பிந்தைய இரு கொள்கைகளும் பின்வருவதைப் போன்ற ஒரு விரிவான கொள்கையிலிருந்து
வருகின்றன என்று முன்வைப்பது நியாயமானது:
% SOURCE CORRECTION TA-STH-019: remove the competing articles from “some a broader principle.”
\begin{enumerate}
	\item[] \stagesinex. முற்றிலும் முடிவிலாப் பல படிநிலைகள் உள்ளன;
	படிநிலைமுறை சாத்தியமான அளவெல்லாம் உயரமானது.
\end{enumerate}
இது ஒரு முறைசாராக் கொள்கை என்பது வெளிப்படை. ஆனால் கணத்தின்
திரள்-படிநிலைமுறைக் கருத்தாக்கத்தால் இது உடனடியாக
\emph{பின்விளைவிக்கப்படாவிட்டாலும்}, அதனுடன் இது நிச்சயமாக
\emph{இசைந்திருப்பதாகத்} தோன்றுகிறது. குறைந்தபட்சம், \limofsize{} போலன்றி,
கணங்கள் படிநிலைப் படிநிலையாக உருவாகின்றன என்ற கருத்தை இது
தக்கவைக்கிறது.

இப்போது \stagesinex{} ஐக் கொண்டு மாற்றீட்டை நியாயப்படுத்துவதே நமது நம்பிக்கை.
அதை எவ்வாறு செய்யலாம் என்று பார்ப்போம்.

\olref[ordinals][idea]{sec} இல், $\omega+\omega$ உடன் (அடிப்படையில்)
ஓரினமாக இருக்க வேண்டிய ஒரு நன்கு வரிசைப்படுத்தலை அமைப்பது எளிது என்று
கண்டோம். வேறு விதமாகக் கூறினால், (அடிப்படையில்) $\omega+\omega$ என்ற
வரிசை வகையைக் கொண்ட படிநிலைப் படிநிலையான ஒரு செயல்முறையை எளிதாகக்
கற்பனை செய்யலாம். எனவே \stagesinex{} ஐ ஏற்றிருந்தால், படிநிலைமுறையின்
குறைந்தபட்சம் ஓர் $\omega+\omega$-ஆம் படிநிலை, அதாவது
$V_{\omega+\omega}$, உள்ளது என்பதையும் நிச்சயமாக ஏற்க வேண்டும்; ஏனெனில்
படிநிலைமுறை இவ்வளவு தூரம் தொடர \emph{முடியும்}.

இச்சிந்தனை பின்வருமாறு பொதுமைப்படுகிறது: எந்த நன்கு வரிசைப்படுத்தலுக்கும்,
படிநிலைமுறையைக் கட்டும் செயல்முறை குறைந்தபட்சம் அந்த நன்கு வரிசைப்படுத்தல்
அளவுக்காவது செல்ல வேண்டும். மேலும்
\olref[ordinals][ordtype]{thmOrdinalRepresentation} ஐ ஓர்
\emph{அடிகோளாகக்} கருதுவதன் மூலம் இதை உறுதிசெய்யலாம். எந்த நன்கு
வரிசைப்படுத்தலும் ஒரு ஃபான் நாய்மன் வரிசையெண்ணுக்கு ஓரினமானது என்று இது
கூறும். ஒவ்வொரு ஃபான் நாய்மன் வரிசையெண்ணும் அதன் சொந்தத் தரநிலைக்குச்
சமமாக இருப்பதால், நமது கணக் கோட்பாட்டில் ஒரு நன்கு வரிசைப்படுத்தலை
விவரிக்க முடியும் போதெல்லாம், கணம் கட்டும் படிநிலைமுறை அந்த நன்கு
வரிசைப்படுத்தலை மீறிச் செல்ல வேண்டும் என்று
\olref[ordinals][ordtype]{thmOrdinalRepresentation} கூறும்.

இக்கருத்து \stagesinex{} இன் ஒரு துணைவிளைவுபோல நிச்சயமாகத் தெரிகிறது.
ஆனால் இக்கருத்திலிருந்து மாற்றீட்டைப் பெறுவதே நமது நோக்கமெனில், எளிய
தொழில்நுட்பத் தடையைச் சந்திக்கிறோம்: மாற்றீடு
\olref[ordinals][ordtype]{thmOrdinalRepresentation} ஐவிடக் கண்டிப்பாக
வலிமையானது. (இக்கவனிப்பை \citet[\S13.2]{Potter2004} முன்வைக்கிறார்;
அதை \olref[sth][replacement][finiteaxiomatizability]{sec} இல் நிறுவுவோம்.)

இதன் விளைவு: மாற்றீட்டைத் தரும் வகையில் \stagesinex{} ஐப் புரிந்துகொள்ள
வேண்டுமெனில், படிநிலைமுறை எந்த நன்கு வரிசைப்படுத்தலையும் மீறிச் செல்கிறது
என்று அது \emph{வெறுமனே} கூற முடியாது. அதைவிட வலிமையான ஒரு கூற்றை அது
முன்வைக்க வேண்டும். இதற்காக \cite{Shoenfield:AST}, இக்கருத்தின் மிகவும்
இயல்பான பின்வரும் வலுவூட்டலை முன்வைத்தார்: படிநிலைமுறை எந்தக் கணத்துடனும்
\emph{இறுதியிணையானதன்று}.\footnote{கோடல் இதைப் போன்ற ஒரு சிந்தனையை
முன்வைத்ததாகத் தெரிகிறது; \citet[p.~223]{Potter2004} ஐப் பார்க்கவும்.
கோடல் மற்றும் \citeauthor{Shoenfield:AST} பற்றிய விவாதத்துக்கு
\citet[90--5]{Incurvati2020} ஐப் பார்க்கவும்.} சற்று விரிவாகக் கூறினால்:
$\tau$ என்பது கணங்களைப் படிநிலைமுறையின் படிநிலைகளுக்கு அனுப்பும் ஒரு
வரைவு எனில், $\tau$-வின்கீழ் எந்தக் கணம் $A$ இன் பிம்பமும் படிநிலைமுறையை
முழுவதுமாக எட்டாது. வேறு விதமாகக் கூறினால் (திட்டவடிவில்):
\begin{enumerate}
	\item[] \stagescofin. $A$ ஒரு கணமாகவும், ஒவ்வொரு $x \in A$ க்கும்
	$\tau(x)$ ஒரு படிநிலையாகவும் இருந்தால், $x \in A$ ஒவ்வொன்றின்
	$\tau(x)$ க்கும் பின்னால் வரும் ஒரு படிநிலை உள்ளது.
\end{enumerate}
\stagescofin{} இன் பொருத்தமாக முறைப்படுத்தப்பட்ட ஒரு வடிவத்தை $\ZF$
நிறுவுகிறது என்பது வெளிப்படை. மறுதிசையில், \stagescofin{} மாற்றீட்டை
நியாயப்படுத்துகிறது என்று முறைசாராமல் வாதிடலாம்.\footnote{\stagescofin{}
இன் ஏதோ ஒரு முறைப்படுத்தலைப் பயன்படுத்தி மாற்றீட்டை நிறுவுவது கடினமாக
இருக்கும்; ஏனெனில் $\Z$ மட்டும் படிநிலைகளை வரையறுக்கும் அளவுக்கு
வலிமையானதன்று, ஆகவே \stagescofin{} ஐ எவ்வாறு முறைப்படுத்துவது என்பது
தெளிவில்லை. எனினும், \olref[sth][replacement][extrinsic]{sec} இல்
விவாதித்தபடி, $\LT$ இன் ஏதோ ஒரு விரிவாக்கத்தில் பணிபுரிவது ஒரு வழி.}
$(\forall x \in A)\lexists![y][\phi(x,y)]$ என்க. ஒவ்வொரு $x \in A$
க்கும் $\sigma(x)$ என்பது $\phi(x,y)$ ஆக அமையும் அந்த $y$ எனவும்,
$\tau(x)$ என்பது $\sigma(x)$ முதன்முதலில் உருவாகும் படிநிலை எனவும் கொள்க.
\stagescofin{} இன்படி, $(\forall x \in A)\tau(x)\in V$ ஆக அமையும் ஒரு
படிநிலை $V$ உள்ளது. இப்போது ஒவ்வொரு $\tau(x) \in V$ ஆகவும்
$\sigma(x) \subseteq \tau(x)$ ஆகவும் இருப்பதால், பிரிப்பைப் பயன்படுத்தி
$\Setabs{y \in V}{(\exists x \in A)\sigma(x) = y} = \Setabs{y}{(\exists x \in
A)\phi(x,y)}$ ஐப் பெறலாம்.

\begin{prob}
	$\ZF$ க்குள் \stagescofin{} ஐ முறைப்படுத்துக.
\end{prob}

ஆகவே \stagescofin{} மாற்றீட்டை உறுதிசெய்கிறது. மேலும் \stagesinex{}
ஆனது \stagescofin{} ஐ உறுதிசெய்கிறது என்பதும் குறைந்தபட்சம் ஏற்புடையதே.
\stagescofin{} தோல்வியுறுகிறது என்க. அப்போது படிநிலைமுறை ஏதோ ஒரு கணம்~$A$
உடன் இறுதியிணையானது; அதாவது, எந்தப் படிநிலை~$S$ க்கும்
$S \in \tau(x)$ ஆக அமையும் ஏதோ ஓர் $x \in A$ இருக்குமாறு ஒரு வரைவு
$\tau$ நமக்குள்ளது. அந்த நிலையில் படிநிலைமுறையின் ``முற்றுமுடிவிலியை''
ஒருவகையில் பிடித்துக்கொள்ள முடிகிறது: $A$ க்குப் பயன்படுத்தப்பட்ட
$\tau$-வின் வீச்சு அதை \emph{முழுவதுமாக எட்டுகிறது}. இது படிநிலைமுறை
``முற்றிலும் முடிவிலியானது'' என்ற சிந்தனையைச் சிதைக்கிறது.
எதிர்மறை வாதமாக: \stagesinex{} ஆனது \stagescofin{} ஐப் பின்விளைவிக்கிறது;
அது மாற்றீட்டை நியாயப்படுத்துகிறது.

மாற்றீட்டுக்கு உள்ளார்ந்த நியாயப்படுத்தலை வழங்கும் உண்மையாகவே நம்பிக்கையூட்டும்
முயற்சி இது. ஆனால் இறுதியில் அது வெற்றிபெறுகிறதா இல்லையா என்பதை நீங்கள்
முடிவுசெய்வதற்கே விட்டுவிட வேண்டும்.

\end{document}
